\documentclass[11pt]{article}
\usepackage[margin=1in]{geometry}
\usepackage{hyperref}
\usepackage{enumitem}
\title{Robot Logic and Conditional Coding}
\author{Piter Garcia}
\date{\today}
\begin{document}
\maketitle

\section*{Grade and Class Match}
\textbf{Grade:} Grade 5\\
\textbf{Likely blocks:} Day 2 · 10:45-11:35 · Pickett; Day 3 · 10:45-11:35 · McGlashon; Day 5 · 10:45-11:35 · Pecora

\section*{Lesson Focus}
Conditional logic, event-based behavior, and partner debugging in robotics.

\section*{Learning Targets}
\begin{itemize}[leftmargin=*]
\item Students use conditional or event-based structures in robot code.
\item Students explain how a program changes based on the condition or trigger.
\item Students debug a challenge with a partner.
\end{itemize}

\section*{Materials}
\begin{itemize}[leftmargin=*]
\item Robots
\item Student devices
\item Challenge instructions
\end{itemize}

\section*{Suggested Lesson Flow}
\begin{enumerate}[leftmargin=*]
\item Open with the exact device, tool, or routine students need in order to start successfully.
\item Model one example task before students begin independent or partner work.
\item Give students time to build, code, design, or document their work while circulating for support.
\item Close with a short share-out, reflection, or debugging discussion.
\end{enumerate}

\section*{Source Notes}
This lesson packet is enriched with transcript-derived lesson notes, the school schedule roster, and official starting-point resources. See the paired Markdown guide for clickable links.

\bibliographystyle{plain}
\bibliography{robot-logic-and-conditional-coding}
\end{document}
